\section{结论}

本文把一部通俗明史著抽取出的知识图谱当作观测对象，检查它的信息结构能否被计算，
以及算出来的结构是否与作者的笔法错位。三个结论：政治攻讦关系占该部全部关系的比重
从第 1--2 部的 6.4\% 升到第 6--7 部的 16.0\%（$2.48$ 倍），而同期的敌对冲突不动（$0.98$ 倍）；
193 条攻讦中 86.0\% 跨阵营，端点集与敌对关系只有 0.219 的 Jaccard，两者盯的不是同一批人；
厂卫系 29 人在图谱里的平均度数 3.76，与全库 3.78 无差别（$0.99$），
而同一套制度的内廷宦官 51 人平均度数 6.10（$1.61$ 倍）——机构隐形，代理人上位。

首条限制写在讨论节：攻讦类别与派系标签与关系抽取共享同一份语料，
「晚期上升」不能排除是抽取顺序的副产物，只有第二份语料的平行抽取才能检验。

下一步可检验的命题是：把同一批 1231 人放进第二部明史著再抽一次，
若攻讦占比的时间形状与该语的七部结构同构，则同源自证可以退居次要，
本篇的三个数字才谈得上从「一本书的图谱」升格为关于明代信息结构的弱证据。
